\section{Trasmissione dell'informazione multimediale}
\subsection{Fasi di trasmissione dell'informazione}
\begin{figure}[htbp]
	\centering
	\begin{tikzpicture}[
		% Stile per i blocchi di elaborazione generici
		block/.style={
			rectangle,
			draw=blue!70!black,
			fill=blue!5,
			thick,
			align=center,
			rounded corners=3pt,
			font=\footnotesize
		},
		% Stile specifico per il canale di trasmissione
		channel/.style={
			rectangle,
			draw=orange!85!black,
			fill=orange!8,
			thick,
			align=center,
			rounded corners=3pt,
			font=\footnotesize
		},
		% Stile per le frecce di collegamento
		arrow/.style={
			-{Stealth[scale=1.2]},
			thick
		},
		% Stile per le etichette dei segnali
		signal/.style={
			align=center
		},
		% Stile per le parentesi
		brace/.style={
			decorate,
			decoration={brace, amplitude=10pt, raise=6pt},
			thick
		}
	]
		% Catena di blocchi
		\node [signal] at (0, 0) (tx) {\textbf{Input}};
		\node [block] at (2.5, 0) (camp) {Sampling};
		\node [block] at (5, 0) (quant) {Quantizer};
		\node [block] at (7.5, 0) (invmap) {Inverse bit\\mapper};
		\node [block] at (10, 0) (p2s) {P/S\\converter};
		\node [block] at (13, 0) (cod) {Codificatore};
		\node [channel] at (13,-1.25) (canale) {Rete o canale\\di trasmissione};
		\node [block] at (13, -2.5) (decod) {Decodificatore};
		\node [block] at (10, -2.5) (s2p) {S/P\\converter};
		\node [block] at (7.5, -2.5) (bmap) {Bit mapper};
		\node [block] at (2.5, -2.5) (interp) {Interpolatore};
		\node [signal] at (0, -2.5) (out) {\textbf{Output}};
		
		% Collegamenti ed etichette
		\draw [arrow] (tx) -- node[signal, below=0.1cm] {$s(t)$ \\ $t \in \mathbb{R}$} (camp);
		\draw [arrow] (camp) -- node[signal, below=0.1cm] {$s_c(n)$ \\ $n \in \mathbb{Z}$} (quant);
		\draw [arrow] (quant) -- node[signal, below=0.1cm] {$s_q(n)$ \\ $n \in \mathbb{Z}$} (invmap);
		\draw [arrow] (invmap) -- node[signal, below=0.1cm] {$\vec{s}_w(n)$ \\ $n \in \mathbb{Z}$} (p2s);
		\draw [arrow] (p2s) -- node[signal, below=0.1cm] {$s_d(n)$ \\ $n \in \mathbb{Z}$} (cod);
		\draw [arrow] (cod) -- (canale);
		\draw [arrow] (canale) -- (decod);
		\draw [arrow] (decod) -- node[signal, above=0.1cm] {$s'_d(n)$ \\ $n \in \mathbb{Z}$} (s2p);
		\draw [arrow] (s2p) -- node[signal, above=0.1cm] {$\vec{s}\,'_w(n)$ \\ $n \in \mathbb{Z}$} (bmap);
		\draw [arrow] (bmap) -- node[signal, above=0.1cm] {$s'_c(n)$ \\ $n \in \mathbb{Z}$} (interp);
		\draw [arrow] (interp) -- node[signal, above=0.1cm] {$s'(t)$ \\ $t \in \mathbb{R}$} (out);

		% --- Parentesi e testi extra ---
		\draw [brace] ($(camp.north west)-(0.2,0)$) -- ($(p2s.north east)+(0.2,-0.17)$)
			node [midway,above=17pt, font=\small] {conversione analogico-digitale - ADC};
		\draw [brace] ($(s2p.south east)+(0.2,0.17)$) -- ($(interp.south west)-(0.2,0)$)
			node [midway, below=17pt, font=\small] {conversione digitale-analogico - DAC};
	\end{tikzpicture}
\end{figure}

\noindent
Il processo di trasmissione dell'informazione di segnali analogici attraverso servizi multimediali digitali prevede le
seguenti fasi:
\begin{enumerate}
	\item conversione analogico-digitale (ADC): conversione dei segnali analogici in segnali digitali, comprende le
	operazioni di campionamento, quantizzazione, e conversione dei simboli in sequenze seriali di bit
	\item codifica (COD): compressione delle sequenze di bit per ridurre la quantità di informazione da trasmettere
	e aggiunta di ridondanza per la protezione dai possibili errori di trasmissione
	\item trasmissione (NET): invio dei simboli attraverso un canale di comunicazione della rete presente
	\item decodifica (DEC): conversione dei simboli ricevuti in segnali digitali
	\item conversione digitale-analogico (DAC): ricostruzione dei segnali analogici a partire dai segnali digitali
	ricevuti, comprende la conversione dei simboli in sequenze seriali di bit e una fase di interpolazione
\end{enumerate}

\subsubsection*{Domini dei segnali analogici e digitali}
Un \textbf{segnale analogico} è una funzione di una o più variabili dipendenti continue in funzione di altre variabili
indipendenti continue.
\[f \; : \; \mathcal{X} \subseteq \mathbb{R}^m \; \to \; \mathcal{Y} \subseteq \mathbb{R}^n \qquad\qquad f(\mathcal{X}, \mathcal{Y}) \; \in \; \mathbb{R}^m \times \mathbb{R}^n\]
\begin{itemize}
	\item segnale audio analogico su nastro magnetico: \(f(t) \in \mathbb{R} \times \mathbb{R}\)
	\item segnale immagine analogico su diapositiva: \(f(x,y) \in \mathbb{R}^2 \times \mathbb{R}\)
	\item segnale video analogico, non campionato: \(f(x,y,t) \in \mathbb{R}^3 \times \mathbb{R}\)
\end{itemize}

\noindent
Un \textbf{segnale analogico campionato} è un segnale analogico in cui le variabili indipendenti sono state discretizzate
attraverso un'acquisizione a intervalli regolari dei valori delle variabili dipendenti. Tale processo di acquisizione è
detto campionamento (approfondito in seguito).
\[f \; : \; \mathcal{X} \subseteq \mathbb{Z}^m \; \to \; \mathcal{Y} \subseteq \mathbb{R}^n \qquad\qquad f(\mathcal{X}, \mathcal{Y}) \; \in \; \mathbb{Z}^m \times \mathbb{R}^n\]
\begin{itemize}
	\item segnale audio campionato: \(f(k) \in \mathbb{Z} \times \mathbb{R}\)
	\item segnale immagine campionato: \(f(m,n) \in \mathbb{Z}^2 \times \mathbb{R}\)
	\item segnale video campionato: \(f(m,n,k) \in \mathbb{Z}^3 \times \mathbb{R}\)
\end{itemize}

\noindent
Un \textbf{segnale digitale} è una funzione di una o più variabili dipendenti discrete in funzione di altre variabili
indipendenti discrete. Si ottiene assegnando un valore discreto alle variabili dipendenti del segnale analogico campionato
attraverso un processo di quantizzazione (approfondito in seguito).
\[f \; : \; \mathcal{X} \subseteq \mathbb{Z}^m \; \to \; \mathcal{Y} \subseteq \mathbb{Z}^n \qquad\qquad f(\mathcal{X}, \mathcal{Y}) \; \in \; \mathbb{Z}^m \times \mathbb{Z}^n\]
\begin{itemize}
	\item segnale audio digitale su CD: \(f(n) \in \mathbb{Z} \times \mathbb{Z}\)
	\item segnale immagine digitale su monitor: \(f(m,n) \in \mathbb{Z}^2 \times \mathbb{Z}\)
	\item segnale video digitale su DVD: \(f(m,n,k) \in \mathbb{Z}^3 \times \mathbb{Z}\)
\end{itemize}

\newpage

\subsection{Campionamento}
\subsubsection*{Trasformata di Fourier}
Dato un segnale a tempo continuo \(s(t)\) (con durata e ampiezza finita), la sua trasformata di Fourier restituisce lo
spettro \(S(f)\) che rappresenta il contenuto frequenziale del segnale \(s(t)\).
\begin{itemize}
	\item la trasformata di Fourier di un tono puro sinusoidale \(s(t) = A \cdot \sin(2\pi f_0 t)\) è un impulso a frequenza
	\(f_0\) con area costante
	\item la trasformata di Fourier del segnale sinc(\(t\)) è un segnale rect(\(f\)), il sinc viene detto segnale a banda
	limitata, proprio perché la sua trasformata di Fourier è limitata a un intervallo di frequenze \([f_m, f_M]\)
\end{itemize}

\subsubsection*{Teorema del campionamento di Shannon e criterio di Nyquist}
Un segnale \(s(t)\) può essere ricostruito con errore nullo a partire dalla successione dei campioni \(s_c(n)\)
ottenuti dal suo campionamento se e solo se tale segnale è banda limitata con frequenza massima \(f_M\) e la
frequenza di campionamento \(F_c\) è almeno il doppio della frequenza massima \(f_M\) del segnale \(s(t)\).
\[s_c(n) = s(t) \big|_{t=n \cdot T_c} \qquad\qquad \text{con} \; F_c = 1 / T_c \geq 2 f_M\]
La relazione \(F_c \geq 2 f_M\) è detta criterio di Nyquist.

\subsubsection*{Formula di interpolazione ideale}
Se il criterio di Nyquist è soddisfatto, il segnale \(s(t)\) può essere ricostruito a partire dai campioni \(s_c(n)\)
attraverso la formula di interpolazione ideale (convoluzione tra segnale campionato e segnale sinc(\(t\))):
\[s(t) = \sum_{n \in \mathbb{Z}} s_c(n) \cdot \text{sinc}\!\left(\frac{t - n T_c}{T_c}\right)\]
La formula di interpolazione ideale consiste nel sommare tanti segnali sinc(\(t\)) centrati negli istanti di campionamento
\(n T_c\), scalati orizzontalmente per occupare tutto il periodo di campionamento \([-T_c, T_c]\) e pesati dai valori dei
campioni ai rispettivi istanti \(s_c(n)\).

\subsubsection*{Filtraggio di segnali limitati nel tempo con banda illimitata}
I segnali multimediali reali hanno durata finita e di conseguenza sono a banda illimitata. Non è possibile, quindi, applicare
strettamente il teorema di Shannon-Nyquist, di conseguenza il segnale campionato non può essere ricostruito con errore nullo.

Si nota, però, che la maggior parte dei segnali si concentra solo su una banda limitata di frequenze. Di conseguenza,
è possibile applicare un filtro passa-basso con frequenza di taglio \(f_c/2\) senza perdere informazioni rilevanti del
segnale originale e senza introdurre errori di distorsione percepibili. In questo modo si limita il contenuto frequenziale
del segnale e risulta possibile applicare il teorema di Shannon-Nyquist garantendo una ricostruzione del segnale campionato
con errore trascurabile.

\subsubsection*{Aliasing - frequenze spurie e sovrapposizione delle repliche dello spettro}
Se si campiona un segnale che non rispetta il criterio di Nyquist, si verifica un fenomeno di aliasing, che consiste nello
spostamento verso il basso delle alte frequenze del segnale campionato che non possono essere ricostruite correttamente.
In altre parole nel segnale ricostruito sono presenti nuove frequenze spurie che non erano presenti nel segnale originale.

Dal punto di vista della trasformata di Fourier, il campionamento di un segnale \(s(t)\) equivale alla replica in
frequenza del suo spettro \(S(f)\) compreso nel range \([-f_M, f_M]\) a intervalli regolari di \(F_c\). Il fenomeno
di aliasing corrisponde alla sovrapposizione (overlapping) di tali repliche se \(f_M > F_c/2\).

Per limitare l'aliasing, si utilizza per l'appunto un filtro passa-basso sul segnale originale prima del campionamento
per ridurre le alte frequenze che generano frequenze spurie o in altre parole per ridurre la banda ed evitare l'overlapping
delle repliche.

\newpage


\subsection{Quantizzazione}
\subsubsection*{Quantizzazione per livelli}
Il processo di quantizzazione consiste nel discretizzare l'ampiezza di un segnale (si suppone già campionato) in un numero
finito di \(L\) livelli. Si definisce, inoltre, la quantità di informazione o costo di codifica \(R = \log_2 L\) come
il numero di bit necessari a rappresentare ogni livello.

Aumentando il numero di livelli \(L\) aumenta la fedeltà del segnale quantizzato rispetto al segnale originale, ma aumenta
anche il costo di codifica \(R\) per rappresentare ogni campione quantizzato.

\subsubsection*{Quantizzazione uniforme}
In un quantizzatore uniforme si suddividere l'intervallo di ampiezza del segnale in \(L\) livelli uniformemente distribuiti,
ciascuno con un passo di quantizzazione \(\Delta\) costante. Il processo di quantizzazione consiste nel rappresentare ogni
campione con il più vicino livello di quantizzazione (o multiplo di \(\Delta\)).
\[Q(x, \Delta) = \Delta \cdot \text{round}\!\left( \frac{x}{\Delta} \right) \qquad \text{con} \; x \in [-A, A], \quad L = 2A / \Delta, \quad \varepsilon = \Delta / 2 \]

\subsubsection*{Quantizzatori mid-rise, mid-tread e dead zone}
In base a come viene trattato il valore zero, i quantizzatori uniformi vengono classificati in:
\begin{itemize}
	\item \textbf{mid-rise}: il valore zero è un valore di soglia tra due livelli di quantizzazione, se il segnale presenta
	piccoli valori di rumore attorno allo zero, questi vengono amplificati e degenerano in uno dei due livelli adiacenti
	\item \textbf{mid-tread}: il valore zero è associato ad un livello di quantizzazione centrato su zero e passo di
	quantizzazione \(\Delta\) uguale a quello degli altri livelli, se il segnale presenta piccoli valori di rumore attorno
	allo zero, questi vengono eliminati e il segnale quantizzato viene riportato a zero
	\item \textbf{dead zone}: come il mid-rise, ma con un passo di quantizzazione maggiore rispetto agli altri livelli,
	serve per comprimere ulteriormente i dati sparsi (con tanti valori piccoli non significativi)
\end{itemize}

\subsubsection*{Rumore di quantizzazione e distorsione}
Il processo di quantizzazione porta inevitabilmente ad una perdita di informazione in quanto è un'operazione irreversibile.
Inoltre si osserva che la qualità del segnale degrada molto velocemente al diminuire del numero di livelli \(L\). Per questo
motivo si utilizzano \(\Delta\) piccoli e \(L\) elevati per garantire una buona fedeltà al segnale originale e non introdurre
distorsioni o artefatti percepibili. Di solito si utilizzano quantizzatori a 8 bit per canale di colore (RGB) e 16 bit per
campionamento audio.

La quantizzazione può essere utilizzata una seconda volta come tecnica di compressione all'interno dei processi di codifica
e può introdurre distorsioni evidenti allo scopo di diminuire la dimensione dei dati. La prima quantizzazione della
conversione DAC, però, non deve indurre distorsioni percepibili.

\subsubsection*{Costo della quantizzazione}
Il costo della quantizzazione è dato dal numero di bit necessari a rappresentare ogni campione quantizzato, ovvero al tasso
di codifica \(R = \log_2 L\) bit per campione.

\subsubsection*{Mean Square Error (MSE) come indice della distorsione}
Si definisce l'errore di quantizzazione \(e(n)\) come la differenza tra il campione del segnale campionato \(s_c(n)\) e
il campione corrispondente del segnale quantizzato \(s_q(n)\):
\[e(n) = s_c(n) - s_q(n)\]

\noindent
Il MSE (Mean Square Error) è un indice della distorsione dovuta alla quantizzazione, corrisponde all'errore quadratico
medio ed equivale alla norma euclidea al quadrato del vettore errore o alla distanza euclidea al quadrato tra il vettore
dei campioni del segnale campionato e il vettore dei campioni del segnale quantizzato (al netto della normalizzazione
per il fattore \(1/N\)).
\[\text{MSE}(s_c(n), s_q(n)) \; = \; \frac{1}{N} \sum_{n=0}^{N-1} [e(n)]^2 \; = \; \frac{1}{N} \sum_{n=0}^{N-1} [s_c(n) - s_q(n)]^2 \; = \; \frac{1}{N} ||\vec{e} \, ||^2 \; = \; \frac{1}{N} ||\vec{s_c} - \vec{s_q}||^2\]

\subsubsection*{Peak Signal-to-Noise Ratio (PSNR) come indice della qualità}
Il PSNR (Peak Signal-to-Noise Ratio) è un indice della qualità del segnale quantizzato, viene calcolato come rapporto
in dB tra la potenza di picco del segnale e il MSE. Il PSNR e il MSE sono strettamente legati tra di loro e totalmente
definiti l'uno dall'altro.
Il PSNR non corrisponde esattamente alla qualità percepita da un osservatore umano, però viene spesso utilizzato data
la sua semplicità di calcolo.
\[\text{PSNR} = 10 \log_{10} \frac{\left( 2^b - 1\right)^2}{\text{MSE}}\]

\noindent
Di seguito è riportato un esempio di calcolo del PSNR per un'immagine in scala di grigi quantizzata a 8 bit con 255
livelli di grigio, con alcune indicazioni sulla qualità percepita in base al valore del PSNR.
\[\text{PSNR} = 10 \log_{10} \frac{255^2}{\text{MSE}} = 48 \text{dB} - 10 \log_{10} \text{MSE}\]
\begin{itemize}
	\item PSNR \(>\) 45 dB: ottima qualità, distorsioni non percepibili
	\item 40 dB \(<\) PSNR \(<\) 45 dB: qualità molto buona, distorsioni difficilmente percepibili
	\item 35 dB \(<\) PSNR \(<\) 40 dB: qualità buona, distorsioni percepibili
	\item 30 dB \(<\) PSNR \(<\) 35 dB: qualità accettabile, distorsioni evidenti
	\item PSNR \(<\) 30 dB: qualità scadente
\end{itemize}

\subsubsection*{Assenza di un blocco inverso alla quantizzazione}
Siccome la quantizzazione è un'operazione irreversibile, non esiste alcun processo inverso in grado di ricostruire
il segnale campionato originale a partire dal segnale quantizzato. Per questo motivo, durante la fase di ricostruzione
del segnale analogico, non ci sarà alcun blocco inverso alla quantizzazione.

\subsection{Inverse bit mapper (IBM) e parallel-to-serial converter (P/S)}
Il quantizzatore produce in output simboli discreti che coincidono con i vari livelli di quantizzazione. Tali simboli
vengono poi convertiti in vettori (o sequenze) di bit tramite il modulo Inverse Bit Mapper (IBM).

L'IBM produce in output vettori di bit in parallelo, per cui è necessario un altro modulo detto Parallel to Serial
Converter (P/S) per convertirli in sequenze seriali che potranno poi essere codificate e trasmesse.

Il bitrate del blocco di conversione ADC è dato dal prodotto tra il tasso di codifica \(R\) del quantizzatore e la
frequenza di campionamento \(F_c\) del segnale campionato.
\(R_\text{ADC} = F_c \cdot R_{Q} = F_c \cdot \log_2 L\)

\subsection{Serial-to-parallel converter (S/P) e bit mapper (BM)}
Una volta che il ricevitore riceve i simboli trasmessi, questi vengono decodificati in sequenze seriali di bit. Tali
sequenze vengono poi convertite in vettori di bit in parallelo tramite il modulo Serial to Parallel Converter (S/P)
per poi essere trasformate in simboli discreti tramite il modulo Bit Mapper (BM). In questo modo si riottiene una
ricostruzione del segnale quantizzato originale che poi verrà interpolato per ottenere il segnale analogico ricostruito.

\subsection{Interpolazione}
\subsubsection*{Formula di interpolazione ideale}
Dal teorema di Shannon-Nyquist, per ricostruire un segnale analogico a partire dai campioni del segnale campionato
secondo opportune condizioni, si utilizza la formula di interpolazione ideale:
\[s(t) = \sum_{n \in \mathbb{Z}} s(n T_c) \cdot \text{sinc}\!\left(\frac{t - n T_c}{T_c}\right)\]
Si nota però che tale formula non è implementabile in pratica, in quanto richiede di sommare un numero infinito di
segnali sinc(\(t\)) a supporto infinito. Nella pratica si sostituisce la sinc(\(t\)) con altri segnali a supporto
finito detti kernel di interpolazione.

\subsubsection*{Kernel di interpolazione e vincoli}
Un kernel di interpolazione è un segnale a supporto finito (implementabile nella pratica) che viene utilizzato al
posto del kernel ideale sinc(\(t\)). La funzione di interpolazione diventa:
\[s(t) = \sum_{n \in \mathbb{Z}} s(n T_c) \cdot h\!\left(\frac{t - n T_c}{T_c}\right) \qquad \text{con} \;\; h\!\left(\frac{t - n T_c}{T_c}\right) \;\; \begin{array}{c}
	\text{kernel di interpolazione dilatato} \\
	\text{nell'intervallo di camp. } [-T_c, T_c] \text{ e}\\
	\text{traslato nell'istante di camp. } n T_c
\end{array}\]
I kernel di interpolazione devono rispettare alcuni vincoli per la conservazione dei valori negli istanti di campionamento
per una corretta ricostruzione del segnale campionato, come illustrato di seguito.
\[h(n) = \begin{cases}
	1 & \text{se } n = 0 \\
	0 & \text{se } n \in \mathbb{Z} \setminus \! \{0\}
\end{cases} \implies \tilde{s}(kT_c) = \sum_{n \in \mathbb{Z}} s(nT_c) \cdot h\!\left(\frac{kT_c - nT_c}{T_c}\right) = \sum_{n \in \mathbb{Z}} s(nT_c) \cdot h(k\!-\!n) = s(kT_c)\]

\subsubsection*{Esempi di kernel di interpolazione}
\begin{itemize}
	\item \textbf{sample and hold}, con kernel di interpolazione a gradino rettangolare:
	\[h(t) = \text{rect}(t-1/2) = \begin{cases}
		1 & \text{se } 0 \leq t < 1 \\
		0 & \text{altrimenti}
	\end{cases}\]
	il segnale ricostruito è un segnale a gradino che mantiene il valore del campione per tutto l'intervallo di campionamento,
	è il kernel di interpolazione più semplice da implementare, ma introduce distorsioni evidenti e alte frequenze spurie
	dovute alla discontinuità delle funzioni di interpolazione, tali artefatti risultano accettabili solo se la \(f_c\) è elevata
	
	\item \textbf{interpolazione lineare} o di primo ordine, con kernel di interpolazione a triangolo:
	\[h(t) = \text{triang}(t) = \begin{cases}
		1 - |t| & \text{se } |t| < 1 \\
		0 & \text{altrimenti}
	\end{cases}\]
	il segnale ricostruito corrisponde ad una linea spezzata che unisce i valori dei vari campioni, è più complesso da
	implementare, ma introduce meno distorsioni e frequenze spurie rispetto al sample and hold
	
	\item \textbf{interpolazione cubica} o di terzo ordine, con kernel di interpolazione polinomiale di terzo grado: \\
	si effettua un'interpolazione con polinomi di terzo grado tra i campioni, è più complesso da implementare e spesso
	si usa negli algoritmi di interpolazione digitale delle immagini (es. cambio di risoluzione di un'immagine) in quanto
	esistono algoritmi efficienti basati sul filtraggio numerico

	\item \textbf{interpolazione con spline cardinali} con kernel di interpolazione polinomiale di grado \(n\): \\
	è possibile usare come kernel di interpolazione un generico polinomio di grado \(n\) che implica la continuità delle
	derivate fino all'ordine \(n-1\) del segnale ricostruito, riduce sempre di più le distorsioni e le frequenze spurie,
	ma risulta sempre più complesso da implementare, per \(n \to \infty\) coincide con l'interpolazione ideale a sinc(\(t\))
	
	\item \textbf{interpolazione con sinc troncato} con kernel di interpolazione a sinc troncato:
	\[h(t) = \text{sinc}(t) \cdot \text{rect}\!\left(\frac{t}{T}\right) = \begin{cases}
		\frac{\sin(\pi t)}{\pi t} & \text{se } |t| < T \\
		0 & \text{altrimenti}
	\end{cases}\]
	è il kernel utilizzato all'interno delle schede audio di fascia alta, in quanto rappresenta un'ottima approssimazione
	fisicamente realizzabile del kernel ideale sinc(\(t\)); il parametro \(T\) viene scelto in fase di progettazione,
	maggiore è \(T\), maggiore è la fedeltà della ricostruzione ed aumenta la complessità computazionale in quanto cresce
	il numero di campioni coinvolti per la ricostruzione del segnale in ogni istante di tempo \(t\)
\end{itemize}
